\chapter{Kết luận và hướng phát triển}

% Phụ trách: ............

Chương này tổng kết những kết quả đạt được trong quá trình thực hiện
đề tài, đánh giá các hạn chế còn tồn tại của hệ thống và đề xuất các
hướng phát triển trong tương lai nhằm nâng cao hiệu quả quản lý
bãi giữ xe.

% ==================================================
\section{Kết quả đạt được}
% ==================================================

Sau thời gian nghiên cứu, phân tích, thiết kế và triển khai, nhóm đã
hoàn thành hệ thống quản lý bãi giữ xe trên nền tảng Windows Forms
và SQL Server.

Hệ thống đáp ứng được các mục tiêu đã đề ra ở Chương~1 và hiện thực
thành công các yêu cầu chức năng được phân tích ở Chương~3.

Các kết quả chính đạt được gồm:

\begin{itemize}
    \item Xây dựng được cơ sở dữ liệu quản lý thông tin loại xe,
    thẻ xe, khách hàng, vé tháng, bảng giá và lịch sử lượt gửi.

    \item Triển khai thành công chức năng ghi nhận xe vào và xe ra
    theo quy trình thực tế của bãi giữ xe.

    \item Xây dựng module tính phí tự động dựa trên loại xe và thời
    gian gửi bằng cách sử dụng \texttt{DateTime}, \texttt{TimeSpan}
    và delegate \texttt{Func}.

    \item Triển khai chức năng quản lý vé tháng, bao gồm đăng ký mới,
    gia hạn và kiểm tra trạng thái còn hiệu lực.

    \item Bảo đảm tính toàn vẹn dữ liệu thông qua transaction trong
    quá trình thanh toán và cập nhật thông tin xe ra.

    \item Hỗ trợ tra cứu lịch sử gửi xe theo biển số, thẻ xe và
    khoảng thời gian.

    \item Hỗ trợ in vé xe khi xe vào và in biên nhận khi thanh toán.

    \item Thống kê doanh thu theo ngày hoặc khoảng thời gian phục vụ
    công tác quản lý và đối soát.

    \item Tổ chức mã nguồn theo mô hình ba lớp (GUI, BLL, DAL),
    giúp chương trình dễ bảo trì và mở rộng.
\end{itemize}

Kết quả kiểm thử ở Chương~6 cho thấy các chức năng chính của hệ thống
hoạt động ổn định, cho kết quả đúng với mong đợi và đáp ứng được các
yêu cầu nghiệp vụ đã đặt ra.

% ==================================================
\section{Hạn chế}
% ==================================================

Mặc dù đã hoàn thành các chức năng cốt lõi của bài toán quản lý
bãi giữ xe, hệ thống vẫn còn một số hạn chế:

\begin{itemize}
    \item Việc nhận dạng biển số xe vẫn được thực hiện bằng cách
    nhập thủ công, chưa hỗ trợ nhận diện tự động bằng camera.

    \item Chưa tích hợp với các thiết bị phần cứng như camera giám
    sát, đầu đọc RFID hoặc barie tự động.

    \item Chưa hỗ trợ thanh toán điện tử bằng ví điện tử, ngân hàng
    hoặc mã QR.

    \item Chưa có cơ chế sao lưu và phục hồi dữ liệu tự động.

    \item Chưa hỗ trợ mô hình quản lý nhiều bãi xe đồng thời trong
    cùng một hệ thống.

    \item Chức năng thống kê còn ở mức cơ bản, chủ yếu tập trung vào
    doanh thu và chưa cung cấp các báo cáo chuyên sâu.

    \item Hệ thống được xây dựng dưới dạng ứng dụng desktop nên chỉ
    hoạt động trên môi trường Windows và chưa hỗ trợ truy cập từ xa
    qua trình duyệt web.
\end{itemize}

Những hạn chế trên chủ yếu xuất phát từ phạm vi của đề tài môn học
và giới hạn về thời gian thực hiện.

% ==================================================
\section{Hướng phát triển}
% ==================================================

Trong tương lai, hệ thống có thể được mở rộng theo các hướng sau:

\subsection{Tích hợp nhận dạng biển số tự động}

Kết hợp camera và các thư viện xử lý ảnh để tự động nhận diện biển
số xe khi xe vào hoặc xe ra. Việc này giúp giảm thao tác nhập liệu
bằng tay và tăng độ chính xác trong quản lý.

\subsection{Hiển thị số chỗ trống}

Hệ thống có thể theo dõi sức chứa của bãi xe và hiển thị số lượng
chỗ còn trống theo thời gian thực. Chức năng này hỗ trợ nhân viên
điều phối xe hiệu quả hơn.

\subsection{Lưu ảnh xe và biển số}

Khi xe vào hoặc xe ra, hệ thống có thể lưu ảnh chụp biển số và
phương tiện để phục vụ đối chiếu khi xảy ra tranh chấp hoặc mất vé.

\subsection{Mở rộng chức năng thống kê}

Bổ sung các báo cáo:

\begin{itemize}
    \item Số lượt xe theo ngày.
    \item Số lượt xe theo loại phương tiện.
    \item Biểu đồ doanh thu theo tháng.
    \item Tỷ lệ sử dụng bãi xe.
    \item Danh sách vé tháng sắp hết hạn.
\end{itemize}

\subsection{Tích hợp thanh toán điện tử}

Cho phép khách hàng thanh toán thông qua mã QR, ví điện tử hoặc
chuyển khoản ngân hàng nhằm giảm thời gian xử lý tại cổng ra.

\subsection{Triển khai trên nền tảng Web}

Chuyển hệ thống từ ứng dụng WinForms sang kiến trúc Web hoặc
Web API giúp người quản lý có thể theo dõi và điều hành bãi xe
từ nhiều thiết bị khác nhau.

\subsection{Quản lý nhiều bãi xe}

Mở rộng hệ thống theo hướng cho phép quản lý đồng thời nhiều
chi nhánh bãi giữ xe với cơ sở dữ liệu tập trung.

% ==================================================
\section*{Kết chương}
% ==================================================

Đề tài ``Quản lý bãi giữ xe'' đã hoàn thành các mục tiêu đặt ra
và xây dựng được một hệ thống đáp ứng các nghiệp vụ cơ bản trong
thực tế như ghi nhận xe vào, xe ra, tính phí, quản lý vé tháng,
tra cứu lịch sử và thống kê doanh thu.

Mặc dù vẫn còn một số hạn chế, hệ thống đã thể hiện được việc
vận dụng các kiến thức về lập trình hướng đối tượng, xử lý thời
gian, cơ sở dữ liệu quan hệ, ADO.NET, transaction và kiến trúc
phân lớp trong quá trình phát triển phần mềm.

Đây là nền tảng để tiếp tục mở rộng và hoàn thiện hệ thống trong
các nghiên cứu và dự án tiếp theo.